% =============================================================================
% CHAPTER 3: EXPERIMENTS AND EVALUATION
% =============================================================================
% COMPLETELY REWRITTEN — Old Knowledge Graph / Data Fabric content removed
% Project: A Context-Aware Framework for Streaming Data Quality Monitoring
% Date: April 24, 2026
% =============================================================================

\section{Introduction to Chapter}
\label{sec:ch3-intro}

This chapter describes the experimental evaluation of StreamDQ. We begin with the
experimental setup (\ref{sec:experimental-setup}): datasets, infrastructure,
and synthetic anomaly injection. \ref{sec:evaluation-methodology} presents the
ground-truth tracking protocol, metrics, ablation design, and statistical tests.
\ref{sec:results-syn-sem} through \ref{sec:results-latency} present benchmark
results as placeholder tables labeled \texttt{[TIER-2\ ESTIMATED]}, pending
implementation and benchmark execution. \ref{sec:ch3-limitations} presents
the mandatory limitations section with all seven documented constraints.

\section{Experimental Setup}
\label{sec:experimental-setup}

\subsection{Datasets}

\textbf{NYC TLC Yellow Taxi.} We use the NYC Taxi and Limousine Commission
Yellow Taxi trip record dataset for January 2024 \citep{nyc_tlc}, stored as
Parquet files with approximately 3 million records. These records contain
zone-level spatial references (PULocationID and DOLocationID, mapping to 263
taxi zones) but \textbf{no GPS latitude/longitude coordinates}. The TLC
dataset is replayed through Apache Kafka at real-time speed using a custom
Kafka producer. Because it lacks GPS, only SYN and SEM rules apply; CRS rules
cannot be evaluated on TLC data.

\textbf{NYC MTA Bus GTFS-realtime.} We use the live NYC MTA Bus GTFS-realtime
public feed \citep{nyc_mta_bus}, which publishes VehiclePosition messages for
approximately 500 buses at approximately 30-second intervals per vehicle.
This feed is publicly accessible and requires no API key. The feed URL,
evaluation timestamp, and dataset version are logged for reproducibility.
Because the MTA Bus feed includes GPS latitude/longitude coordinates, all
three rule layers (SYN, SEM, and CRS) apply to this stream.

Figure \ref{fig:experimental-pipeline} shows the evaluation pipeline.

\begin{figure}[htbp]
  \centering
  \fbox{\parbox{0.92\linewidth}{
    \centering
    \textbf{[Figure F7: Experimental Pipeline — see
    \texttt{figures/F7\_Experimental\_Pipeline.excalidraw.md}]}
  }}
  \caption{Evaluation pipeline for StreamDQ. Synthetic anomalies are injected
    at controlled rates into the event stream before rule evaluation. Every
    injected anomaly is tracked in the ground\_truth\_events table. Every
    detected violation is matched against ground truth to compute precision
    and recall.}
  \label{fig:experimental-pipeline}
\end{figure}

\subsection{Infrastructure}

StreamDQ operates in two execution modes:

\textbf{LocalPipeline (development mode).} An in-process Python pipeline
processes events from Kafka or a replay log. State is maintained in Python
dictionaries (in-memory). SQLite is used for violation storage. LocalPipeline
is used for unit tests, development, and functional validation.
\texttt{[TIER-2\ ESTIMATED]}

\textbf{FlinkPipeline (distributed mode).} Apache Flink processes events
from Kafka in a distributed cluster. State is maintained in RocksDB.
PostgreSQL stores violations and ground-truth data. FlinkPipeline is used
for full benchmark runs.
\texttt{[TIER-2\ ESTIMATED]}

\textbf{Critical distinction.} LocalPipeline and FlinkPipeline have
different latency and throughput characteristics. LocalPipeline results are
not equivalent to distributed Flink results. All latency numbers reported
from LocalPipeline are labeled ``LocalPipeline profiling only, not
equivalent to distributed Flink performance.''

Figure \ref{fig:dashboard-mockup} shows the monitoring dashboard.

\begin{figure}[htbp]
  \centering
  \fbox{\parbox{0.92\linewidth}{
    \centering
    \textbf{[Figure F8: Grafana Dashboard Mockup — see
    \texttt{figures/F8\_Dashboard\_Mockup.excalidraw.md}]}
  }}
  \caption{Grafana dashboard for StreamDQ monitoring. Four panels display:
    (1) violation rate by rule, (2) TQS composite score over time,
    (3) L0 coverage fraction, and (4) processing latency percentiles.
    Screenshot placeholder: replace with actual dashboard after deployment.}
  \label{fig:dashboard-mockup}
\end{figure}

\subsection{Synthetic Anomaly Injection}

We inject synthetic anomalies into the event stream to create ground truth
for evaluation. Each injected anomaly is tagged with its entity index,
anomaly type, and injection rate. Table \ref{tab:anomaly-types} summarizes
the anomaly types and injection parameters.

\begin{table}[htbp]
  \centering
  \caption{Synthetic anomaly types and injection parameters.
    NG-4: CRS003 recall is \texttt{[TIER-3\ UNMEASURABLE]}
    because the replay suppression gate blocks synthetic duplicate injection.}
  \label{tab:anomaly-types}
  \begin{tabular}{lllr}
    \toprule
    Anomaly Type & Rule Triggered & Injection Method & Status \\
    \midrule
    NULL field         & SYN001 & Set required field to NULL       & Ready \\
    NaN field          & SYN001 & Set float field to NaN          & Ready \\
    Negative fare      & SYN002 & Set fare\_amount $<$ 0           & Ready \\
    GPS speed spike   & CRS001 & Two positions 2 km apart in 30 s & Ready \\
    GPS jump          & CRS002 & Positions $>$400 m at $\leq$30 s & Ready \\
    Duplicate event   & CRS003 & Emit original + duplicate         & NG-4 blocked \\
    \bottomrule
  \end{tabular}
\end{table}

Injection rates are 0.5\%, 1.0\%, 2.0\%, and 5.0\% of events per anomaly type.
Each trial uses a different random seed (SEED = 42, 43, 44). A warmup period
of 10,000 events (approximately 60 seconds at real-time speed) precedes each
measurement window. The measurement window is 600 seconds (10 minutes) per trial.
Three independent trials are run per condition.

\textbf{NG-4 (CRS003 blocker).} The replay suppression gate in the current
code tags synthetic duplicate events as \texttt{is\_replay=True}, which triggers
the gate in the deduplication logic and suppresses violation emission. To measure
CRS003 recall, synthetic duplicate events must be tagged \texttt{is\_replay=False}.
This is an evaluation infrastructure bug, not a rule logic bug. CRS003 recall
is labeled \texttt{[TIER-3\ UNMEASURABLE]} until NG-4 is fixed.

\section{Evaluation Methodology}
\label{sec:evaluation-methodology}

\subsection{Ground-Truth Tracking}

We track every injected anomaly in a PostgreSQL \texttt{ground\_truth\_events}
table with the following schema:

\begin{itemize}
  \itemsep0em
  \item \texttt{run\_id}: UUID identifying the evaluation run
  \item \texttt{injection\_rate}: float (0.005, 0.01, 0.02, 0.05)
  \item \texttt{anomaly\_type}: VARCHAR (NULL, NAN, NEGATIVE\_FARE, GPS\_SPEED,
        GPS\_JUMP, DUPLICATE)
  \item \texttt{entity\_index}: INTEGER (position in event stream)
  \item \texttt{injected\_value}: JSONB (field name, injected value)
  \item \texttt{injected\_at}: TIMESTAMP
\end{itemize}

Every detected violation is matched against ground truth using exact
entity index matching. No fuzzy matching is used.

\subsection{Metrics}

Table \ref{tab:eval-metrics} presents the evaluation metrics.

\begin{table}[htbp]
  \centering
  \caption{Evaluation metrics, confidence interval methods, and bootstrap
    specifications. Bootstrap iterations should be increased to 10,000 for
    publication.}
  \label{tab:eval-metrics}
  \begin{tabular}{lllr}
    \toprule
    Metric & CI Method & Bootstrap Iterations & Formula \\
    \midrule
    Precision, Recall, F1 & Wilson score     & $\geq$ 1,000 & $2P\!R/(P+R)$ \\
    Violation rate        & Clopper-Pearson & $\geq$ 1,000 & $n_{\text{viol}}/N$ \\
    P99 latency           & Percentile       & $\geq$ 1,000 & $P_{99}(L)$ \\
    Throughput            & t-test bootstrap & $\geq$ 1,000 & events/s \\
    TQS composite         & Percentile       & $\geq$ 1,000 & Weighted sum \\
    Spearman $\rho$       & Fisher z-transform& $\geq$ 1,000 & Rank corr. \\
    \bottomrule
  \end{tabular}
\end{table}

All confidence intervals are at the 95\% level using the percentile bootstrap
method (1,000 iterations). For publication, bootstrap iterations should be
increased to 10,000.
\texttt{[TODO: increase bootstrap to 10,000 iterations for publication]}

\textbf{Latency measurement.} Processing latency is measured from event arrival
(at Kafka produce timestamp) to violation stored (PostgreSQL INSERT complete).
This is an end-to-end measurement, not an internal processing time.
Latency numbers from LocalPipeline are labeled as such.

\textbf{TQS correlation.} TQS composite scores are correlated with downstream
ETA prediction error, not with injection rate. ETA prediction error is
computed independently from GTFS static schedule data and published real-time
delays from the MTA API. This design avoids the circularity of measuring
TQS against injection rate.

\subsection{Ablation Studies}

We conduct three ablation studies to answer the research questions:

\begin{itemize}
  \itemsep0em
  \item \textbf{RQ1 ablation}: Static (global L4) thresholds vs. L0--L4
        context-aware thresholds. Measures whether finer context improves
        detection F1 at L0 cells.
  \item \textbf{RQ4 ablation}: SYN-only vs. SYN+SEM vs. SYN+SEM+CRS.
        Measures whether each rule layer adds incremental F1.
  \item \textbf{RQ6 ablation}: Rule-only vs. rule+ML (Bayesian Optimization only).
        Measures whether ML calibration improves F1.
\end{itemize}

\subsection{Statistical Tests}

Table \ref{tab:statistical-tests} presents the statistical tests.

\begin{table}[htbp]
  \centering
  \caption{Statistical tests, significance levels, and corrections.
    Holm--Bonferroni is applied within test families.}
  \label{tab:statistical-tests}
  \begin{tabular}{lllr}
    \toprule
    Test & Purpose & $\alpha$ & Correction \\
    \midrule
    Wilcoxon signed-rank & Non-parametric paired comparison & 0.05 & Holm--Bonferroni \\
    McNemar's mid-p exact & Paired binary classifier comparison & 0.01 & Holm--Bonferroni \\
    Friedman's test & Multiple related samples (L0--L4) & 0.05 & Nemenyi post-hoc \\
    Bootstrap percentile & Non-parametric CI & 0.05 & Per-method \\
    \bottomrule
  \end{tabular}
\end{table}

Holm--Bonferroni correction is applied within test families at $\alpha = 0.05$.
For RQ1 (5 pairwise comparisons), $\alpha_{\text{adj}} = 0.05 / 5 = 0.01$.
For RQ2 (6 pairwise comparisons), $\alpha_{\text{adj}} = 0.05 / 6 \approx 0.0083$.

\section{Results: SYN/SEM Rules}
\label{sec:results-syn-sem}

Table \ref{tab:results-syn-sem} presents precision, recall, and F1 for SYN
and SEM rules on the NYC TLC Yellow Taxi dataset.

\begin{table}[htbp]
  \centering
  \caption{Precision, recall, and F1 for SYN and SEM rules on NYC TLC.
    Results are \texttt{[TIER-2\ ESTIMATED]} pending benchmark execution.
    All values are accompanied by 95\% bootstrap confidence intervals
    (1,000 iterations). TODO: increase bootstrap to 10,000 iterations.}
  \label{tab:results-syn-sem}
  \begin{tabular}{lcccccc}
    \toprule
    \multirow{2}{*}{Rule} & \multirow{2}{*}{Precision} & \multirow{2}{*}{Recall}
    & \multirow{2}{*}{F1} & \multirow{2}{*}{95\% CI} &
    \multirow{2}{*}{n$_{\text{inj}}$} \\
    & & & & & \\
    \midrule
    SYN001 (NULL)      & [TIER-2\ ESTIMATED] & [TIER-2\ ESTIMATED]
                       & [TIER-2\ ESTIMATED] & — & — \\[4pt]
    SYN001 (NaN)       & [TIER-2\ ESTIMATED] & [TIER-2\ ESTIMATED]
                       & [TIER-2\ ESTIMATED] & — & — \\[4pt]
    SYN002 (neg fare)   & [TIER-2\ ESTIMATED] & [TIER-2\ ESTIMATED]
                       & [TIER-2\ ESTIMATED] & — & — \\[4pt]
    SYN003 (bounds)    & [TIER-2\ ESTIMATED] & [TIER-2\ ESTIMATED]
                       & [TIER-2\ ESTIMATED] & — & — \\[4pt]
    SEM001 (fare)      & [TIER-2\ ESTIMATED] & [TIER-2\ ESTIMATED]
                       & [TIER-2\ ESTIMATED] & — & — \\[4pt]
    SEM002 (distance)  & [TIER-2\ ESTIMATED] & [TIER-2\ ESTIMATED]
                       & [TIER-2\ ESTIMATED] & — & — \\[4pt]
    SEM003 (pax count) & [TIER-2\ ESTIMATED] & [TIER-2\ ESTIMATED]
                       & [TIER-2\ ESTIMATED] & — & — \\[4pt]
    GTFSSem002 (stale) & [TIER-2\ ESTIMATED] & [TIER-2\ ESTIMATED]
                       & [TIER-2\ ESTIMATED] & — & — \\
    \bottomrule
  \end{tabular}
\end{table}

The SYN001 (NULL) and SYN001 (NaN) rules provide the baseline precision
for record-level validation. We expect high precision (close to 1.0) for
NULL checks because synthetic NULL injection produces unambiguous violations.
SEM001 and SEM002 are expected to have lower recall because context-aware
thresholds may not fire on injected anomalies that fall within the rolling
percentile bounds.

\section{Results: CRS Rules}
\label{sec:results-crs}

Table \ref{tab:results-crs} presents precision, recall, and F1 for CRS
rules on the NYC MTA Bus GTFS-realtime stream.

\begin{table}[htbp]
  \centering
  \caption{Precision, recall, and F1 for CRS rules on NYC MTA Bus GTFS-realtime.
    CRS001 and CRS002 are \texttt{[TIER-2\ ESTIMATED]} pending benchmark.
    CRS003 is \texttt{[TIER-3\ UNMEASURABLE]} due to NG-4
    (replay suppression gate blocks synthetic duplicate injection).
    All values are accompanied by 95\% bootstrap confidence intervals
    (1,000 iterations). TODO: increase bootstrap to 10,000 iterations.}
  \label{tab:results-crs}
  \begin{tabular}{lcccccc}
    \toprule
    \multirow{2}{*}{Rule} & \multirow{2}{*}{Precision} & \multirow{2}{*}{Recall}
    & \multirow{2}{*}{F1} & \multirow{2}{*}{95\% CI} &
    \multirow{2}{*}{n$_{\text{inj}}$} \\
    & & & & & \\
    \midrule
    CRS001 (GPS speed)   & [TIER-2\ ESTIMATED] & [TIER-2\ ESTIMATED]
                         & [TIER-2\ ESTIMATED] & — & — \\[4pt]
    CRS002 (GPS jump)    & [TIER-2\ ESTIMATED] & [TIER-2\ ESTIMATED]
                         & [TIER-2\ ESTIMATED] & — & — \\[4pt]
    CRS003 (deduplication)& [TIER-2\ ESTIMATED] & \cellcolor[gray]{0.85}
                         \texttt{[TIER-3]} & \cellcolor[gray]{0.85}
                         \texttt{[TIER-3]} & \cellcolor[gray]{0.85} NG-4
                         & \cellcolor[gray]{0.85} — \\
    \bottomrule
  \end{tabular}
\end{table}

CRS001 and CRS002 use physics-based thresholds (speed bounds, jump distance)
that are independent of the data distribution, so we expect high precision.
CRS001 lower bound (2.0 km/h) may produce false positives for stationary
vehicles with GPS jitter; the sustained-low-speed guard ($\Delta t > 60$ s)
mitigates this. CRS002 precision depends on the quality of the synthetic
GPS jump injection; the 400-meter threshold was calibrated for the 30-second
GTFS-realtime update interval.

CRS003 recall is \texttt{[TIER-3\ UNMEASURABLE]} because the replay
suppression gate (NG-4) blocks synthetic duplicate injection. Fixing NG-4
requires tagging synthetic duplicate events as \texttt{is\_replay=False}.

\section{Results: Context-Aware Threshold Ablation}
\label{sec:results-threshold-abl}

Table \ref{tab:results-threshold-abl} presents the ablation results for
context-aware thresholds (RQ1).

\begin{table}[htbp]
  \centering
  \caption{Ablation: static (L4 global) vs. context-aware (L0--L4) thresholds.
    ``$\Delta$F1 at L0 cells'' measures the F1 improvement specifically for
    events falling into L0 context cells.
    ``Aggregate $\Delta$F1'' measures improvement across all events.
    L0 coverage is 2--5\%; aggregate F1 improvement is expected at 1--2 pp.
    Results are \texttt{[TIER-2\ ESTIMATED]} pending benchmark execution.}
  \label{tab:results-threshold-abl}
  \begin{tabular}{lcccccc}
    \toprule
    Configuration & $\Delta$F1 at L0 cells (pp) & 95\% CI
    & Aggregate $\Delta$F1 (pp) & 95\% CI & L0 coverage \\
    \midrule
    Static (L4 global)   & baseline & — & baseline & — & — \\
    L0 only              & [TIER-2] & — & [TIER-2] & — & 2--5\% \\
    L0--L4 (hierarchical)& [TIER-2] & — & [TIER-2] & — & 2--5\% \\
    \bottomrule
  \end{tabular}
\end{table}

\textbf{Critical caveat.} The 5 pp target for $\Delta$F1 applies specifically
to events falling into L0 context cells (2--5\% of events). Because most
events fall to L3 or L4 (where thresholds are similar to static), the
aggregate $\Delta$F1 across all events is expected to be 1--2 pp, not 5 pp.
Both metrics will be reported separately. The aggregate number is likely
underpowered at the current sample size (approximately 90 measurement windows
$\times$ 3 trials).

Statistical significance is assessed via Wilcoxon signed-rank test with
Bonferroni correction ($\alpha_{\text{adj}} = 0.01$ for 5 comparisons).

\section{Results: ML Augmentation}
\label{sec:results-ml}

Table \ref{tab:results-ml} presents the ablation results for ML-augmented
threshold calibration (RQ6).

\begin{table}[htbp]
  \centering
  \caption{Ablation: rule-only vs. rule+Bayesian Optimization.
    Bayesian Optimization runs hourly, tuning k\_multiplier, IF\_alpha,
    and weekend\_discount on F1 objective (not violation rate).
    Isolation Forest is conditional (GO if $\rho_{\text{IF,P90}} < 0.8$).
    XGBoost is scaffold only (training target undefined).
    LSTM is NO-GO (GPS archive unconfirmed; high CRS002 redundancy).
    Results are \texttt{[TIER-2\ ESTIMATED]} pending benchmark execution.
    Positive, negative, and mixed results are all valid findings.}
  \label{tab:results-ml}
  \begin{tabular}{lcccccc}
    \toprule
    Configuration & Precision & Recall & F1 & 95\% CI & n$_{\text{inj}}$ \\
    \midrule
    Rule-only (baseline) & [TIER-2] & [TIER-2] & [TIER-2] & — & — \\
    + Bayesian Optimization & [TIER-2] & [TIER-2] & [TIER-2] & — & — \\
    + Isolation Forest (cond.) & [TIER-2] & [TIER-2] & [TIER-2] & — & — \\
    \bottomrule
  \end{tabular}
\end{table}

This is a \textit{measured contribution}: if ML provides no improvement,
the negative result is scientifically valuable and will be reported as such.

\section{Latency and Throughput}
\label{sec:results-latency}

Table \ref{tab:results-latency} presents latency and throughput measurements.

\begin{table}[htbp]
  \centering
  \caption{Latency and throughput measurements. LocalPipeline results are
    \texttt{[TIER-2\ ESTIMATED]} and labeled ``LocalPipeline profiling only,
    not equivalent to distributed Flink performance.'' Distributed Flink
    benchmarks require infrastructure deployment.}
  \label{tab:results-latency}
  \begin{tabular}{lcccc}
    \toprule
    Mode & P50 latency (ms) & P99 latency (ms) & 95\% CI (ms) &
    Throughput (events/s) \\
    \midrule
    LocalPipeline & [TIER-2] & [TIER-2] & — & [TIER-2] \\
    FlinkPipeline & [TIER-2] & [TIER-2] & — & [TIER-2] \\
    \bottomrule
  \end{tabular}
\end{table}

\textbf{Critical limitation.} LocalPipeline and FlinkPipeline have fundamentally
different architectures. LocalPipeline processes events in-process with Python
state management; FlinkPipeline distributes events across JVM TaskManagers with
RocksDB state. P99 latency from LocalPipeline is not representative of
distributed Flink latency. Distributed benchmarks require Flink cluster
deployment and are not yet available.

\section{Limitations}
\label{sec:ch3-limitations}

We explicitly acknowledge seven limitations of this work.

\textbf{1. CRS003 recall is unmeasurable (NG-4).} The replay suppression
gate tags synthetic duplicate events as \texttt{is\_replay=True}, which
triggers the gate and suppresses violation emission. CRS003 recall
cannot be measured until synthetic duplicates are tagged
\texttt{is\_replay=False}. This is an evaluation infrastructure bug, not
a rule logic bug. CRS003 precision is measurable; recall is labeled
\texttt{[TIER-3\ UNMEASURABLE]}.
\texttt{[TIER-3\ UNMEASURABLE]}

\textbf{2. L0 coverage is sparse (2--5\% of events).} The aggregate
$\Delta$F1 from context-aware thresholds across all events is likely
1--2 pp, not the 5 pp target which applies only to L0 cells. Both
L0-specific and aggregate metrics will be reported separately. The aggregate
metric is likely underpowered at approximately 90 measurement windows.

\textbf{3. ML augmentation results are estimated (TIER-2).} All ML
numbers (F1 improvement from Bayesian Optimization, Isolation Forest
conditional activation) require Phase 3 implementation and benchmark
execution. These results are placeholders.

\textbf{4. Single dataset (NYC TLC + NYC MTA Bus).} Results may not
generalize to other transit agencies, geographic regions, or data schemas.
CRS rules (CRS001, CRS002) calibrated for NYC MTA Bus may need parameter
adjustment for different vehicle types, update frequencies, or geographic
contexts.

\textbf{5. LocalPipeline $\neq$ distributed Flink.} All latency and
throughput numbers from LocalPipeline are labeled as development-mode
profiling only. Distributed Flink benchmarks require cluster deployment
and are not yet available. Claiming Flink-level latency from LocalPipeline
results would be architecturally incorrect.

\textbf{6. L4 global fallback masks context-specific patterns.} Events
evaluated at L4 use the global threshold regardless of their context.
If the global threshold is lenient, L4 events that would violate a
context-specific threshold will not be flagged. This is an inherent
trade-off of the hierarchical fallback design: it ensures coverage but
dilutes context specificity.

\textbf{7. D4 External context is a stub.} The holiday indicator in
D4 (External context dimension) is not implemented. The five-dimensional
context decomposition effectively operates on four dimensions (D1, D2, D3,
D5). Holidays, traffic incidents, and other external signals are not
captured by the current context model.

\section{Chapter Summary}
\label{sec:ch3-summary}

This chapter presented the evaluation methodology and results. We described
the NYC TLC and NYC MTA Bus datasets, the LocalPipeline and FlinkPipeline
infrastructure, and the synthetic anomaly injection protocol with five anomaly
types. We detailed ground-truth tracking, metrics with 95\% bootstrap CI,
ablation studies for RQ1, RQ4, and RQ6, and statistical tests with
Holm--Bonferroni correction. Result tables are placeholders labeled
\texttt{[TIER-2\ ESTIMATED]} pending benchmark execution; CRS003 is
\texttt{[TIER-3\ UNMEASURABLE]} due to NG-4. The mandatory limitations
section explicitly acknowledged all seven documented constraints.
